% A Didacta document.
%
% Two lines of setup, then the composition. Compare the legacy master: 60 lines
% of preamble, four \import calls with `../../../../` paths, a documentclass
% incantation, and the theorem template chosen by hand per language.
%
% Compile it directly and you get the notes in the course language. Pass a
% profile and you get any other output:
%
%   pdflatex "\def\DidactaProfile{slides}% A Didacta document.
%
% Two lines of setup, then the composition. Compare the legacy master: 60 lines
% of preamble, four \import calls with `../../../../` paths, a documentclass
% incantation, and the theorem template chosen by hand per language.
%
% Compile it directly and you get the notes in the course language. Pass a
% profile and you get any other output:
%
%   pdflatex "\def\DidactaProfile{slides}% A Didacta document.
%
% Two lines of setup, then the composition. Compare the legacy master: 60 lines
% of preamble, four \import calls with `../../../../` paths, a documentclass
% incantation, and the theorem template chosen by hand per language.
%
% Compile it directly and you get the notes in the course language. Pass a
% profile and you get any other output:
%
%   pdflatex "\def\DidactaProfile{slides}% A Didacta document.
%
% Two lines of setup, then the composition. Compare the legacy master: 60 lines
% of preamble, four \import calls with `../../../../` paths, a documentclass
% incantation, and the theorem template chosen by hand per language.
%
% Compile it directly and you get the notes in the course language. Pass a
% profile and you get any other output:
%
%   pdflatex "\def\DidactaProfile{slides}\input{tema-1}"
%   didacta build tema-1 --profile slides --language va

\input{didacta-bootstrap}
\usepackage{didacta}

%% Course metadata and the document title come from course.yaml and
%% year.yaml, in whichever language is being built. Didacta injects them.
%%
%% The declarations below are only a fallback, so that opening this file in an
%% editor and pressing compile still produces a titled document. Anything
%% Didacta injects overrides them.
\DidactaCourse{
  title   = {Anàlisi Matemàtica III},
  teacher = {Javier Falcó},
}
\DidactaDocument{Preliminars: espais normats}

\begin{document}
\DidactaTitlePage
\DidactaContents

\section{Espais normats}
\DidactaUnit{analysis/normed-spaces/definition}
\DidactaUnit{analysis/normed-spaces/induced-metric}

\end{document}
"
%   didacta build tema-1 --profile slides --language va

%%
%% didacta-bootstrap.tex -- choose the document class from the output profile.
%%
%% A Didacta document begins with
%%
%%     \input{didacta-bootstrap}
%%     \usepackage{didacta}
%%
%% and nothing else before \begin{document}. The class is not written in the
%% source: it comes from the profile, so one file produces every output.
%%
%% Why a bootstrap file rather than a document class of its own: slides need
%% `beamer` and prose needs `article` or `book`, and no single class can be
%% both. Selecting the class before \documentclass runs is the only way to keep
%% one source. The legacy system reached the same conclusion; Didacta makes the
%% choice explicit and named instead of inferring it from which class happens
%% to be loaded.
%%
%% Two knobs, both settable from the command line:
%%
%%     \def\DidactaProfile{slides}     which output (see didacta-profiles.tex)
%%     \def\DidactaLanguage{va}        es | va | en
%%
%% Both have defaults, so opening the file in an editor and pressing compile
%% produces the written notes in Castilian without any setup. That matters:
%% direct compilation with working SyncTeX is a hard requirement, and it must
%% not depend on remembering to pass flags.
%%
\makeatletter

\@ifundefined{DidactaProfile}{\def\DidactaProfile{notes}}{}
\@ifundefined{DidactaLanguage}{\def\DidactaLanguage{es}}{}

%% Extra class options a document may add, e.g. a paper size.
\@ifundefined{DidactaClassOptions}{\def\DidactaClassOptions{}}{}

%% Let a document override the class outright. Rarely wanted, but an exam in
%% the `exam` class or a thesis in `memoir` should not need a fork of Didacta.
\@ifundefined{DidactaClass}{\def\DidactaClass{}}{}

\input{didacta-profiles}

%% Fail loudly on an unknown profile. A silent fallback would hand back a PDF
%% that looks plausible and is the wrong output.
\@ifundefined{didacta@profile@\DidactaProfile @class}{%
  \begingroup
    \newlinechar=`\^^J
    \errmessage{^^J%
      ! Didacta: unknown output profile `\DidactaProfile'.^^J%
      ! Known profiles:\didacta@profile@list ^^J%
      ! Declare new ones with \string\DidactaDeclareProfile\space in
      didacta-profiles.tex.^^J}%
  \endgroup
}{}

\edef\didacta@class{\@nameuse{didacta@profile@\DidactaProfile @class}}
\edef\didacta@classoptions{\@nameuse{didacta@profile@\DidactaProfile @options}}

%% A document-level class override wins.
\ifx\DidactaClass\@empty\else
  \edef\didacta@class{\DidactaClass}%
\fi

%% Append the document's own options after the profile's, so the document wins
%% on any conflict.
\ifx\DidactaClassOptions\@empty\else
  \edef\didacta@classoptions{\didacta@classoptions,\DidactaClassOptions}%
\fi

%% `notes=show` has to reach beamer as a *class* option, not a later setting,
%% so it is the one axis the bootstrap has to look at. `\in@` is the kernel's
%% substring test, which keeps this dependency-free.
\edef\didacta@axes{\@nameuse{didacta@profile@\DidactaProfile @axes}}
\expandafter\in@\expandafter{\expandafter n\expandafter o\expandafter t%
  \expandafter e\expandafter s\expandafter =\expandafter s\expandafter h%
  \expandafter o\expandafter w\expandafter}\expandafter{\didacta@axes}
\ifin@
  \edef\didacta@classoptions{\didacta@classoptions,notes=show}%
\fi

\expandafter\documentclass\expandafter[\didacta@classoptions]{\didacta@class}

\makeatother

\usepackage{didacta}

%% Course metadata and the document title come from course.yaml and
%% year.yaml, in whichever language is being built. Didacta injects them.
%%
%% The declarations below are only a fallback, so that opening this file in an
%% editor and pressing compile still produces a titled document. Anything
%% Didacta injects overrides them.
\DidactaCourse{
  title   = {Anàlisi Matemàtica III},
  teacher = {Javier Falcó},
}
\DidactaDocument{Preliminars: espais normats}

\begin{document}
\DidactaTitlePage
\DidactaContents

\section{Espais normats}
\DidactaUnit{analysis/normed-spaces/definition}
\DidactaUnit{analysis/normed-spaces/induced-metric}

\end{document}
"
%   didacta build tema-1 --profile slides --language va

%%
%% didacta-bootstrap.tex -- choose the document class from the output profile.
%%
%% A Didacta document begins with
%%
%%     %%
%% didacta-bootstrap.tex -- choose the document class from the output profile.
%%
%% A Didacta document begins with
%%
%%     \input{didacta-bootstrap}
%%     \usepackage{didacta}
%%
%% and nothing else before \begin{document}. The class is not written in the
%% source: it comes from the profile, so one file produces every output.
%%
%% Why a bootstrap file rather than a document class of its own: slides need
%% `beamer` and prose needs `article` or `book`, and no single class can be
%% both. Selecting the class before \documentclass runs is the only way to keep
%% one source. The legacy system reached the same conclusion; Didacta makes the
%% choice explicit and named instead of inferring it from which class happens
%% to be loaded.
%%
%% Two knobs, both settable from the command line:
%%
%%     \def\DidactaProfile{slides}     which output (see didacta-profiles.tex)
%%     \def\DidactaLanguage{va}        es | va | en
%%
%% Both have defaults, so opening the file in an editor and pressing compile
%% produces the written notes in Castilian without any setup. That matters:
%% direct compilation with working SyncTeX is a hard requirement, and it must
%% not depend on remembering to pass flags.
%%
\makeatletter

\@ifundefined{DidactaProfile}{\def\DidactaProfile{notes}}{}
\@ifundefined{DidactaLanguage}{\def\DidactaLanguage{es}}{}

%% Extra class options a document may add, e.g. a paper size.
\@ifundefined{DidactaClassOptions}{\def\DidactaClassOptions{}}{}

%% Let a document override the class outright. Rarely wanted, but an exam in
%% the `exam` class or a thesis in `memoir` should not need a fork of Didacta.
\@ifundefined{DidactaClass}{\def\DidactaClass{}}{}

\input{didacta-profiles}

%% Fail loudly on an unknown profile. A silent fallback would hand back a PDF
%% that looks plausible and is the wrong output.
\@ifundefined{didacta@profile@\DidactaProfile @class}{%
  \begingroup
    \newlinechar=`\^^J
    \errmessage{^^J%
      ! Didacta: unknown output profile `\DidactaProfile'.^^J%
      ! Known profiles:\didacta@profile@list ^^J%
      ! Declare new ones with \string\DidactaDeclareProfile\space in
      didacta-profiles.tex.^^J}%
  \endgroup
}{}

\edef\didacta@class{\@nameuse{didacta@profile@\DidactaProfile @class}}
\edef\didacta@classoptions{\@nameuse{didacta@profile@\DidactaProfile @options}}

%% A document-level class override wins.
\ifx\DidactaClass\@empty\else
  \edef\didacta@class{\DidactaClass}%
\fi

%% Append the document's own options after the profile's, so the document wins
%% on any conflict.
\ifx\DidactaClassOptions\@empty\else
  \edef\didacta@classoptions{\didacta@classoptions,\DidactaClassOptions}%
\fi

%% `notes=show` has to reach beamer as a *class* option, not a later setting,
%% so it is the one axis the bootstrap has to look at. `\in@` is the kernel's
%% substring test, which keeps this dependency-free.
\edef\didacta@axes{\@nameuse{didacta@profile@\DidactaProfile @axes}}
\expandafter\in@\expandafter{\expandafter n\expandafter o\expandafter t%
  \expandafter e\expandafter s\expandafter =\expandafter s\expandafter h%
  \expandafter o\expandafter w\expandafter}\expandafter{\didacta@axes}
\ifin@
  \edef\didacta@classoptions{\didacta@classoptions,notes=show}%
\fi

\expandafter\documentclass\expandafter[\didacta@classoptions]{\didacta@class}

\makeatother

%%     \usepackage{didacta}
%%
%% and nothing else before \begin{document}. The class is not written in the
%% source: it comes from the profile, so one file produces every output.
%%
%% Why a bootstrap file rather than a document class of its own: slides need
%% `beamer` and prose needs `article` or `book`, and no single class can be
%% both. Selecting the class before \documentclass runs is the only way to keep
%% one source. The legacy system reached the same conclusion; Didacta makes the
%% choice explicit and named instead of inferring it from which class happens
%% to be loaded.
%%
%% Two knobs, both settable from the command line:
%%
%%     \def\DidactaProfile{slides}     which output (see didacta-profiles.tex)
%%     \def\DidactaLanguage{va}        es | va | en
%%
%% Both have defaults, so opening the file in an editor and pressing compile
%% produces the written notes in Castilian without any setup. That matters:
%% direct compilation with working SyncTeX is a hard requirement, and it must
%% not depend on remembering to pass flags.
%%
\makeatletter

\@ifundefined{DidactaProfile}{\def\DidactaProfile{notes}}{}
\@ifundefined{DidactaLanguage}{\def\DidactaLanguage{es}}{}

%% Extra class options a document may add, e.g. a paper size.
\@ifundefined{DidactaClassOptions}{\def\DidactaClassOptions{}}{}

%% Let a document override the class outright. Rarely wanted, but an exam in
%% the `exam` class or a thesis in `memoir` should not need a fork of Didacta.
\@ifundefined{DidactaClass}{\def\DidactaClass{}}{}

%%
%% didacta-profiles.tex -- the output profile registry.
%%
%% This file is the single source of truth for what an output *is*. It is read
%% by didacta-bootstrap.tex (which needs the document class) and again by
%% didacta.sty (which needs the behaviour axes), so a profile is described in
%% exactly one place.
%%
%% A profile is a named preset over five independent axes:
%%
%%   medium     slides | document
%%              slides -> beamer;  document -> article/book + beamerarticle,
%%              so that \begin{frame} and \frametitle keep working in prose.
%%
%%   detail     brief | full
%%              brief keeps only what fits on a slide; full adds the prose,
%%              the proofs and the worked examples.
%%
%%   audience   student | teacher
%%              teacher adds presenter notes and teaching guidance.
%%
%%   solutions  hidden | answers | full
%%              answers is the short result; full is the worked solution.
%%
%%   pauses     on | off
%%              off collapses \pause so every slide prints once -- what you
%%              want for a handout or a projector-free reading copy.
%%
%% Adding an output means adding one \DidactaDeclareProfile line. Nothing else
%% in the system needs to know about it.
%%
%% Usage from the command line (this is how the engine drives it):
%%
%%   pdflatex "\def\DidactaProfile{slides}\def\DidactaLanguage{va}\input{master}"
%%
%% This file is \input from two places with different ambient catcodes: the
%% bootstrap (before \documentclass, where @ is `other') and didacta.sty
%% (where @ is already a letter). So it saves and restores the catcode of @
%% instead of ending with a blanket \makeatother, which would silently break
%% whichever caller had it as a letter.
\expandafter\edef\csname didacta@profiles@restorecat\endcsname{%
  \catcode`\noexpand\@=\the\catcode`\@\relax}
\makeatletter

%% Read twice on purpose, so guard the declarations. Re-reading must be a
%% no-op rather than an error.
\@ifundefined{didacta@profiles@loaded}{}{%
  \didacta@profiles@restorecat
  \endinput}
\def\didacta@profiles@loaded{1}

%% \DidactaDeclareProfile{id}{class}{class options}{axes}
%%
%% `axes` is a comma list of `key=value` from the five axes above. Anything
%% unset falls back to the defaults declared further down.
%% Un id que ya está declarado se **sustituye**, y no se vuelve a listar. Es
%% lo que permite que una plantilla del repositorio --que se declara después,
%% desde el fichero que genera el motor-- reemplace a una de las que trae
%% Didacta sin que la lista de perfiles conocidos salga con el nombre dos
%% veces. Sin plantillas esto no cambia nada: ningún id se repite aquí.
\newcommand{\DidactaDeclareProfile}[4]{%
  \@ifundefined{didacta@profile@#1@class}{%
    \expandafter\g@addto@macro\expandafter\didacta@profile@list
      \expandafter{,#1}%
  }{}%
  \@namedef{didacta@profile@#1@class}{#2}%
  \@namedef{didacta@profile@#1@options}{#3}%
  \@namedef{didacta@profile@#1@axes}{#4}%
}

\def\didacta@profile@list{}

%% ---------------------------------------------------------------------------
%% Slide outputs
%% ---------------------------------------------------------------------------

%% What you project in class. Pauses reveal content step by step.
\DidactaDeclareProfile{slides}{beamer}{10pt,notheorems}{%
  medium=slides,detail=brief,audience=student,solutions=hidden,pauses=on}

%% The same slides with every overlay collapsed: one page per slide. This is
%% the copy you hand out or read on a train.
\DidactaDeclareProfile{slides-flat}{beamer}{10pt,handout,notheorems}{%
  medium=slides,detail=brief,audience=student,solutions=hidden,pauses=off}

%% Slides plus the presenter notes and the answers. `notes=show` puts beamer's
%% \note content on facing pages.
\DidactaDeclareProfile{slides-teacher}{beamer}{10pt,handout,notheorems}{%
  medium=slides,detail=brief,audience=teacher,solutions=full,pauses=off,notes=show}

%% ---------------------------------------------------------------------------
%% Document outputs
%% ---------------------------------------------------------------------------

%% The full written notes: everything on the slides, plus the prose and proofs
%% that do not fit on one. This is the student's book.
\DidactaDeclareProfile{notes}{article}{12pt,oneside}{%
  medium=document,detail=full,audience=student,solutions=hidden,pauses=off}

%% The same notes with worked solutions included.
\DidactaDeclareProfile{notes-solutions}{article}{12pt,oneside}{%
  medium=document,detail=full,audience=student,solutions=full,pauses=off}

%% The teacher's copy: full detail, teaching guidance, every solution.
\DidactaDeclareProfile{notes-teacher}{article}{12pt,oneside}{%
  medium=document,detail=full,audience=teacher,solutions=full,pauses=off}

%% A bound volume for a whole course: chapters, two-sided, real front matter.
\DidactaDeclareProfile{book}{book}{11pt,twoside,openright}{%
  medium=document,detail=full,audience=student,solutions=hidden,pauses=off}

%% ---------------------------------------------------------------------------
%% Single-sheet outputs
%% ---------------------------------------------------------------------------

%% A one- or two-page reference sheet, and the format a lab or seminar script
%% is built in. Tight margins, sections numbered but not subsections.
%%
%% Three levels, the same three a problem sheet has: a script with exercises in
%% it is handed out, checked and marked exactly like one.
\DidactaDeclareProfile{handout}{article}{11pt,oneside}{%
  medium=document,detail=full,audience=student,solutions=hidden,pauses=off,
  layout=compact}

%% The same sheet with the result under each exercise, so they can check
%% themselves.
\DidactaDeclareProfile{handout-answers}{article}{11pt,oneside}{%
  medium=document,detail=full,audience=student,solutions=answers,pauses=off,
  layout=compact}

%% The teacher's copy: results, worked solutions and the marking notes.
\DidactaDeclareProfile{handout-teacher}{article}{11pt,oneside}{%
  medium=document,detail=full,audience=teacher,solutions=full,pauses=off,
  layout=compact}

%% ---------------------------------------------------------------------------
%% Problem-sheet outputs
%%
%% The same source produces the three, and they are the three levels of the
%% problem editor: statement, result, worked solution. Level n shows fields 1
%% to n, which is what makes the menu answerable without reading a table.
%%
%% This is the axis the legacy system drove with \def\CurrentAudience and
%% multiaudience; here it is just `solutions`.
%%
%% There is deliberately no fourth `problems-solutions` -- a student copy with
%% the full working and no marking notes. It produced the same PDF as the
%% teacher's copy for every sheet in the repository (nothing uses `marking`),
%% so it was two menu entries for one output, and choosing between them was a
%% guess. The copy with everything in it is the teacher's.
%% ---------------------------------------------------------------------------

%% What the students receive: statements only.
\DidactaDeclareProfile{problems}{article}{12pt,oneside}{%
  medium=document,detail=full,audience=student,solutions=hidden,pauses=off,
  layout=compact}

%% Statements with the short result, so they can check their own work.
\DidactaDeclareProfile{problems-answers}{article}{12pt,oneside}{%
  medium=document,detail=full,audience=student,solutions=answers,pauses=off,
  layout=compact}

%% The teacher's copy: the worked solution step by step, the marking notes and
%% the common mistakes.
\DidactaDeclareProfile{problems-teacher}{article}{12pt,oneside}{%
  medium=document,detail=full,audience=teacher,solutions=full,pauses=off,
  layout=compact}

%% ---------------------------------------------------------------------------
%% Assessment output
%% ---------------------------------------------------------------------------

%% An exam paper. Same content, but numbered space for answers and no hints.
\DidactaDeclareProfile{exam}{article}{12pt,oneside}{%
  medium=document,detail=brief,audience=student,solutions=hidden,pauses=off,
  layout=exam}

%% The marking scheme for that exam.
\DidactaDeclareProfile{exam-marking}{article}{12pt,oneside}{%
  medium=document,detail=full,audience=teacher,solutions=full,pauses=off,
  layout=exam}

\didacta@profiles@restorecat


%% Fail loudly on an unknown profile. A silent fallback would hand back a PDF
%% that looks plausible and is the wrong output.
\@ifundefined{didacta@profile@\DidactaProfile @class}{%
  \begingroup
    \newlinechar=`\^^J
    \errmessage{^^J%
      ! Didacta: unknown output profile `\DidactaProfile'.^^J%
      ! Known profiles:\didacta@profile@list ^^J%
      ! Declare new ones with \string\DidactaDeclareProfile\space in
      didacta-profiles.tex.^^J}%
  \endgroup
}{}

\edef\didacta@class{\@nameuse{didacta@profile@\DidactaProfile @class}}
\edef\didacta@classoptions{\@nameuse{didacta@profile@\DidactaProfile @options}}

%% A document-level class override wins.
\ifx\DidactaClass\@empty\else
  \edef\didacta@class{\DidactaClass}%
\fi

%% Append the document's own options after the profile's, so the document wins
%% on any conflict.
\ifx\DidactaClassOptions\@empty\else
  \edef\didacta@classoptions{\didacta@classoptions,\DidactaClassOptions}%
\fi

%% `notes=show` has to reach beamer as a *class* option, not a later setting,
%% so it is the one axis the bootstrap has to look at. `\in@` is the kernel's
%% substring test, which keeps this dependency-free.
\edef\didacta@axes{\@nameuse{didacta@profile@\DidactaProfile @axes}}
\expandafter\in@\expandafter{\expandafter n\expandafter o\expandafter t%
  \expandafter e\expandafter s\expandafter =\expandafter s\expandafter h%
  \expandafter o\expandafter w\expandafter}\expandafter{\didacta@axes}
\ifin@
  \edef\didacta@classoptions{\didacta@classoptions,notes=show}%
\fi

\expandafter\documentclass\expandafter[\didacta@classoptions]{\didacta@class}

\makeatother

\usepackage{didacta}

%% Course metadata and the document title come from course.yaml and
%% year.yaml, in whichever language is being built. Didacta injects them.
%%
%% The declarations below are only a fallback, so that opening this file in an
%% editor and pressing compile still produces a titled document. Anything
%% Didacta injects overrides them.
\DidactaCourse{
  title   = {Anàlisi Matemàtica III},
  teacher = {Javier Falcó},
}
\DidactaDocument{Preliminars: espais normats}

\begin{document}
\DidactaTitlePage
\DidactaContents

\section{Espais normats}
\DidactaUnit{analysis/normed-spaces/definition}
\DidactaUnit{analysis/normed-spaces/induced-metric}

\end{document}
"
%   didacta build tema-1 --profile slides --language va

%%
%% didacta-bootstrap.tex -- choose the document class from the output profile.
%%
%% A Didacta document begins with
%%
%%     %%
%% didacta-bootstrap.tex -- choose the document class from the output profile.
%%
%% A Didacta document begins with
%%
%%     %%
%% didacta-bootstrap.tex -- choose the document class from the output profile.
%%
%% A Didacta document begins with
%%
%%     \input{didacta-bootstrap}
%%     \usepackage{didacta}
%%
%% and nothing else before \begin{document}. The class is not written in the
%% source: it comes from the profile, so one file produces every output.
%%
%% Why a bootstrap file rather than a document class of its own: slides need
%% `beamer` and prose needs `article` or `book`, and no single class can be
%% both. Selecting the class before \documentclass runs is the only way to keep
%% one source. The legacy system reached the same conclusion; Didacta makes the
%% choice explicit and named instead of inferring it from which class happens
%% to be loaded.
%%
%% Two knobs, both settable from the command line:
%%
%%     \def\DidactaProfile{slides}     which output (see didacta-profiles.tex)
%%     \def\DidactaLanguage{va}        es | va | en
%%
%% Both have defaults, so opening the file in an editor and pressing compile
%% produces the written notes in Castilian without any setup. That matters:
%% direct compilation with working SyncTeX is a hard requirement, and it must
%% not depend on remembering to pass flags.
%%
\makeatletter

\@ifundefined{DidactaProfile}{\def\DidactaProfile{notes}}{}
\@ifundefined{DidactaLanguage}{\def\DidactaLanguage{es}}{}

%% Extra class options a document may add, e.g. a paper size.
\@ifundefined{DidactaClassOptions}{\def\DidactaClassOptions{}}{}

%% Let a document override the class outright. Rarely wanted, but an exam in
%% the `exam` class or a thesis in `memoir` should not need a fork of Didacta.
\@ifundefined{DidactaClass}{\def\DidactaClass{}}{}

\input{didacta-profiles}

%% Fail loudly on an unknown profile. A silent fallback would hand back a PDF
%% that looks plausible and is the wrong output.
\@ifundefined{didacta@profile@\DidactaProfile @class}{%
  \begingroup
    \newlinechar=`\^^J
    \errmessage{^^J%
      ! Didacta: unknown output profile `\DidactaProfile'.^^J%
      ! Known profiles:\didacta@profile@list ^^J%
      ! Declare new ones with \string\DidactaDeclareProfile\space in
      didacta-profiles.tex.^^J}%
  \endgroup
}{}

\edef\didacta@class{\@nameuse{didacta@profile@\DidactaProfile @class}}
\edef\didacta@classoptions{\@nameuse{didacta@profile@\DidactaProfile @options}}

%% A document-level class override wins.
\ifx\DidactaClass\@empty\else
  \edef\didacta@class{\DidactaClass}%
\fi

%% Append the document's own options after the profile's, so the document wins
%% on any conflict.
\ifx\DidactaClassOptions\@empty\else
  \edef\didacta@classoptions{\didacta@classoptions,\DidactaClassOptions}%
\fi

%% `notes=show` has to reach beamer as a *class* option, not a later setting,
%% so it is the one axis the bootstrap has to look at. `\in@` is the kernel's
%% substring test, which keeps this dependency-free.
\edef\didacta@axes{\@nameuse{didacta@profile@\DidactaProfile @axes}}
\expandafter\in@\expandafter{\expandafter n\expandafter o\expandafter t%
  \expandafter e\expandafter s\expandafter =\expandafter s\expandafter h%
  \expandafter o\expandafter w\expandafter}\expandafter{\didacta@axes}
\ifin@
  \edef\didacta@classoptions{\didacta@classoptions,notes=show}%
\fi

\expandafter\documentclass\expandafter[\didacta@classoptions]{\didacta@class}

\makeatother

%%     \usepackage{didacta}
%%
%% and nothing else before \begin{document}. The class is not written in the
%% source: it comes from the profile, so one file produces every output.
%%
%% Why a bootstrap file rather than a document class of its own: slides need
%% `beamer` and prose needs `article` or `book`, and no single class can be
%% both. Selecting the class before \documentclass runs is the only way to keep
%% one source. The legacy system reached the same conclusion; Didacta makes the
%% choice explicit and named instead of inferring it from which class happens
%% to be loaded.
%%
%% Two knobs, both settable from the command line:
%%
%%     \def\DidactaProfile{slides}     which output (see didacta-profiles.tex)
%%     \def\DidactaLanguage{va}        es | va | en
%%
%% Both have defaults, so opening the file in an editor and pressing compile
%% produces the written notes in Castilian without any setup. That matters:
%% direct compilation with working SyncTeX is a hard requirement, and it must
%% not depend on remembering to pass flags.
%%
\makeatletter

\@ifundefined{DidactaProfile}{\def\DidactaProfile{notes}}{}
\@ifundefined{DidactaLanguage}{\def\DidactaLanguage{es}}{}

%% Extra class options a document may add, e.g. a paper size.
\@ifundefined{DidactaClassOptions}{\def\DidactaClassOptions{}}{}

%% Let a document override the class outright. Rarely wanted, but an exam in
%% the `exam` class or a thesis in `memoir` should not need a fork of Didacta.
\@ifundefined{DidactaClass}{\def\DidactaClass{}}{}

%%
%% didacta-profiles.tex -- the output profile registry.
%%
%% This file is the single source of truth for what an output *is*. It is read
%% by didacta-bootstrap.tex (which needs the document class) and again by
%% didacta.sty (which needs the behaviour axes), so a profile is described in
%% exactly one place.
%%
%% A profile is a named preset over five independent axes:
%%
%%   medium     slides | document
%%              slides -> beamer;  document -> article/book + beamerarticle,
%%              so that \begin{frame} and \frametitle keep working in prose.
%%
%%   detail     brief | full
%%              brief keeps only what fits on a slide; full adds the prose,
%%              the proofs and the worked examples.
%%
%%   audience   student | teacher
%%              teacher adds presenter notes and teaching guidance.
%%
%%   solutions  hidden | answers | full
%%              answers is the short result; full is the worked solution.
%%
%%   pauses     on | off
%%              off collapses \pause so every slide prints once -- what you
%%              want for a handout or a projector-free reading copy.
%%
%% Adding an output means adding one \DidactaDeclareProfile line. Nothing else
%% in the system needs to know about it.
%%
%% Usage from the command line (this is how the engine drives it):
%%
%%   pdflatex "\def\DidactaProfile{slides}\def\DidactaLanguage{va}\input{master}"
%%
%% This file is \input from two places with different ambient catcodes: the
%% bootstrap (before \documentclass, where @ is `other') and didacta.sty
%% (where @ is already a letter). So it saves and restores the catcode of @
%% instead of ending with a blanket \makeatother, which would silently break
%% whichever caller had it as a letter.
\expandafter\edef\csname didacta@profiles@restorecat\endcsname{%
  \catcode`\noexpand\@=\the\catcode`\@\relax}
\makeatletter

%% Read twice on purpose, so guard the declarations. Re-reading must be a
%% no-op rather than an error.
\@ifundefined{didacta@profiles@loaded}{}{%
  \didacta@profiles@restorecat
  \endinput}
\def\didacta@profiles@loaded{1}

%% \DidactaDeclareProfile{id}{class}{class options}{axes}
%%
%% `axes` is a comma list of `key=value` from the five axes above. Anything
%% unset falls back to the defaults declared further down.
%% Un id que ya está declarado se **sustituye**, y no se vuelve a listar. Es
%% lo que permite que una plantilla del repositorio --que se declara después,
%% desde el fichero que genera el motor-- reemplace a una de las que trae
%% Didacta sin que la lista de perfiles conocidos salga con el nombre dos
%% veces. Sin plantillas esto no cambia nada: ningún id se repite aquí.
\newcommand{\DidactaDeclareProfile}[4]{%
  \@ifundefined{didacta@profile@#1@class}{%
    \expandafter\g@addto@macro\expandafter\didacta@profile@list
      \expandafter{,#1}%
  }{}%
  \@namedef{didacta@profile@#1@class}{#2}%
  \@namedef{didacta@profile@#1@options}{#3}%
  \@namedef{didacta@profile@#1@axes}{#4}%
}

\def\didacta@profile@list{}

%% ---------------------------------------------------------------------------
%% Slide outputs
%% ---------------------------------------------------------------------------

%% What you project in class. Pauses reveal content step by step.
\DidactaDeclareProfile{slides}{beamer}{10pt,notheorems}{%
  medium=slides,detail=brief,audience=student,solutions=hidden,pauses=on}

%% The same slides with every overlay collapsed: one page per slide. This is
%% the copy you hand out or read on a train.
\DidactaDeclareProfile{slides-flat}{beamer}{10pt,handout,notheorems}{%
  medium=slides,detail=brief,audience=student,solutions=hidden,pauses=off}

%% Slides plus the presenter notes and the answers. `notes=show` puts beamer's
%% \note content on facing pages.
\DidactaDeclareProfile{slides-teacher}{beamer}{10pt,handout,notheorems}{%
  medium=slides,detail=brief,audience=teacher,solutions=full,pauses=off,notes=show}

%% ---------------------------------------------------------------------------
%% Document outputs
%% ---------------------------------------------------------------------------

%% The full written notes: everything on the slides, plus the prose and proofs
%% that do not fit on one. This is the student's book.
\DidactaDeclareProfile{notes}{article}{12pt,oneside}{%
  medium=document,detail=full,audience=student,solutions=hidden,pauses=off}

%% The same notes with worked solutions included.
\DidactaDeclareProfile{notes-solutions}{article}{12pt,oneside}{%
  medium=document,detail=full,audience=student,solutions=full,pauses=off}

%% The teacher's copy: full detail, teaching guidance, every solution.
\DidactaDeclareProfile{notes-teacher}{article}{12pt,oneside}{%
  medium=document,detail=full,audience=teacher,solutions=full,pauses=off}

%% A bound volume for a whole course: chapters, two-sided, real front matter.
\DidactaDeclareProfile{book}{book}{11pt,twoside,openright}{%
  medium=document,detail=full,audience=student,solutions=hidden,pauses=off}

%% ---------------------------------------------------------------------------
%% Single-sheet outputs
%% ---------------------------------------------------------------------------

%% A one- or two-page reference sheet, and the format a lab or seminar script
%% is built in. Tight margins, sections numbered but not subsections.
%%
%% Three levels, the same three a problem sheet has: a script with exercises in
%% it is handed out, checked and marked exactly like one.
\DidactaDeclareProfile{handout}{article}{11pt,oneside}{%
  medium=document,detail=full,audience=student,solutions=hidden,pauses=off,
  layout=compact}

%% The same sheet with the result under each exercise, so they can check
%% themselves.
\DidactaDeclareProfile{handout-answers}{article}{11pt,oneside}{%
  medium=document,detail=full,audience=student,solutions=answers,pauses=off,
  layout=compact}

%% The teacher's copy: results, worked solutions and the marking notes.
\DidactaDeclareProfile{handout-teacher}{article}{11pt,oneside}{%
  medium=document,detail=full,audience=teacher,solutions=full,pauses=off,
  layout=compact}

%% ---------------------------------------------------------------------------
%% Problem-sheet outputs
%%
%% The same source produces the three, and they are the three levels of the
%% problem editor: statement, result, worked solution. Level n shows fields 1
%% to n, which is what makes the menu answerable without reading a table.
%%
%% This is the axis the legacy system drove with \def\CurrentAudience and
%% multiaudience; here it is just `solutions`.
%%
%% There is deliberately no fourth `problems-solutions` -- a student copy with
%% the full working and no marking notes. It produced the same PDF as the
%% teacher's copy for every sheet in the repository (nothing uses `marking`),
%% so it was two menu entries for one output, and choosing between them was a
%% guess. The copy with everything in it is the teacher's.
%% ---------------------------------------------------------------------------

%% What the students receive: statements only.
\DidactaDeclareProfile{problems}{article}{12pt,oneside}{%
  medium=document,detail=full,audience=student,solutions=hidden,pauses=off,
  layout=compact}

%% Statements with the short result, so they can check their own work.
\DidactaDeclareProfile{problems-answers}{article}{12pt,oneside}{%
  medium=document,detail=full,audience=student,solutions=answers,pauses=off,
  layout=compact}

%% The teacher's copy: the worked solution step by step, the marking notes and
%% the common mistakes.
\DidactaDeclareProfile{problems-teacher}{article}{12pt,oneside}{%
  medium=document,detail=full,audience=teacher,solutions=full,pauses=off,
  layout=compact}

%% ---------------------------------------------------------------------------
%% Assessment output
%% ---------------------------------------------------------------------------

%% An exam paper. Same content, but numbered space for answers and no hints.
\DidactaDeclareProfile{exam}{article}{12pt,oneside}{%
  medium=document,detail=brief,audience=student,solutions=hidden,pauses=off,
  layout=exam}

%% The marking scheme for that exam.
\DidactaDeclareProfile{exam-marking}{article}{12pt,oneside}{%
  medium=document,detail=full,audience=teacher,solutions=full,pauses=off,
  layout=exam}

\didacta@profiles@restorecat


%% Fail loudly on an unknown profile. A silent fallback would hand back a PDF
%% that looks plausible and is the wrong output.
\@ifundefined{didacta@profile@\DidactaProfile @class}{%
  \begingroup
    \newlinechar=`\^^J
    \errmessage{^^J%
      ! Didacta: unknown output profile `\DidactaProfile'.^^J%
      ! Known profiles:\didacta@profile@list ^^J%
      ! Declare new ones with \string\DidactaDeclareProfile\space in
      didacta-profiles.tex.^^J}%
  \endgroup
}{}

\edef\didacta@class{\@nameuse{didacta@profile@\DidactaProfile @class}}
\edef\didacta@classoptions{\@nameuse{didacta@profile@\DidactaProfile @options}}

%% A document-level class override wins.
\ifx\DidactaClass\@empty\else
  \edef\didacta@class{\DidactaClass}%
\fi

%% Append the document's own options after the profile's, so the document wins
%% on any conflict.
\ifx\DidactaClassOptions\@empty\else
  \edef\didacta@classoptions{\didacta@classoptions,\DidactaClassOptions}%
\fi

%% `notes=show` has to reach beamer as a *class* option, not a later setting,
%% so it is the one axis the bootstrap has to look at. `\in@` is the kernel's
%% substring test, which keeps this dependency-free.
\edef\didacta@axes{\@nameuse{didacta@profile@\DidactaProfile @axes}}
\expandafter\in@\expandafter{\expandafter n\expandafter o\expandafter t%
  \expandafter e\expandafter s\expandafter =\expandafter s\expandafter h%
  \expandafter o\expandafter w\expandafter}\expandafter{\didacta@axes}
\ifin@
  \edef\didacta@classoptions{\didacta@classoptions,notes=show}%
\fi

\expandafter\documentclass\expandafter[\didacta@classoptions]{\didacta@class}

\makeatother

%%     \usepackage{didacta}
%%
%% and nothing else before \begin{document}. The class is not written in the
%% source: it comes from the profile, so one file produces every output.
%%
%% Why a bootstrap file rather than a document class of its own: slides need
%% `beamer` and prose needs `article` or `book`, and no single class can be
%% both. Selecting the class before \documentclass runs is the only way to keep
%% one source. The legacy system reached the same conclusion; Didacta makes the
%% choice explicit and named instead of inferring it from which class happens
%% to be loaded.
%%
%% Two knobs, both settable from the command line:
%%
%%     \def\DidactaProfile{slides}     which output (see didacta-profiles.tex)
%%     \def\DidactaLanguage{va}        es | va | en
%%
%% Both have defaults, so opening the file in an editor and pressing compile
%% produces the written notes in Castilian without any setup. That matters:
%% direct compilation with working SyncTeX is a hard requirement, and it must
%% not depend on remembering to pass flags.
%%
\makeatletter

\@ifundefined{DidactaProfile}{\def\DidactaProfile{notes}}{}
\@ifundefined{DidactaLanguage}{\def\DidactaLanguage{es}}{}

%% Extra class options a document may add, e.g. a paper size.
\@ifundefined{DidactaClassOptions}{\def\DidactaClassOptions{}}{}

%% Let a document override the class outright. Rarely wanted, but an exam in
%% the `exam` class or a thesis in `memoir` should not need a fork of Didacta.
\@ifundefined{DidactaClass}{\def\DidactaClass{}}{}

%%
%% didacta-profiles.tex -- the output profile registry.
%%
%% This file is the single source of truth for what an output *is*. It is read
%% by didacta-bootstrap.tex (which needs the document class) and again by
%% didacta.sty (which needs the behaviour axes), so a profile is described in
%% exactly one place.
%%
%% A profile is a named preset over five independent axes:
%%
%%   medium     slides | document
%%              slides -> beamer;  document -> article/book + beamerarticle,
%%              so that \begin{frame} and \frametitle keep working in prose.
%%
%%   detail     brief | full
%%              brief keeps only what fits on a slide; full adds the prose,
%%              the proofs and the worked examples.
%%
%%   audience   student | teacher
%%              teacher adds presenter notes and teaching guidance.
%%
%%   solutions  hidden | answers | full
%%              answers is the short result; full is the worked solution.
%%
%%   pauses     on | off
%%              off collapses \pause so every slide prints once -- what you
%%              want for a handout or a projector-free reading copy.
%%
%% Adding an output means adding one \DidactaDeclareProfile line. Nothing else
%% in the system needs to know about it.
%%
%% Usage from the command line (this is how the engine drives it):
%%
%%   pdflatex "\def\DidactaProfile{slides}\def\DidactaLanguage{va}\input{master}"
%%
%% This file is \input from two places with different ambient catcodes: the
%% bootstrap (before \documentclass, where @ is `other') and didacta.sty
%% (where @ is already a letter). So it saves and restores the catcode of @
%% instead of ending with a blanket \makeatother, which would silently break
%% whichever caller had it as a letter.
\expandafter\edef\csname didacta@profiles@restorecat\endcsname{%
  \catcode`\noexpand\@=\the\catcode`\@\relax}
\makeatletter

%% Read twice on purpose, so guard the declarations. Re-reading must be a
%% no-op rather than an error.
\@ifundefined{didacta@profiles@loaded}{}{%
  \didacta@profiles@restorecat
  \endinput}
\def\didacta@profiles@loaded{1}

%% \DidactaDeclareProfile{id}{class}{class options}{axes}
%%
%% `axes` is a comma list of `key=value` from the five axes above. Anything
%% unset falls back to the defaults declared further down.
%% Un id que ya está declarado se **sustituye**, y no se vuelve a listar. Es
%% lo que permite que una plantilla del repositorio --que se declara después,
%% desde el fichero que genera el motor-- reemplace a una de las que trae
%% Didacta sin que la lista de perfiles conocidos salga con el nombre dos
%% veces. Sin plantillas esto no cambia nada: ningún id se repite aquí.
\newcommand{\DidactaDeclareProfile}[4]{%
  \@ifundefined{didacta@profile@#1@class}{%
    \expandafter\g@addto@macro\expandafter\didacta@profile@list
      \expandafter{,#1}%
  }{}%
  \@namedef{didacta@profile@#1@class}{#2}%
  \@namedef{didacta@profile@#1@options}{#3}%
  \@namedef{didacta@profile@#1@axes}{#4}%
}

\def\didacta@profile@list{}

%% ---------------------------------------------------------------------------
%% Slide outputs
%% ---------------------------------------------------------------------------

%% What you project in class. Pauses reveal content step by step.
\DidactaDeclareProfile{slides}{beamer}{10pt,notheorems}{%
  medium=slides,detail=brief,audience=student,solutions=hidden,pauses=on}

%% The same slides with every overlay collapsed: one page per slide. This is
%% the copy you hand out or read on a train.
\DidactaDeclareProfile{slides-flat}{beamer}{10pt,handout,notheorems}{%
  medium=slides,detail=brief,audience=student,solutions=hidden,pauses=off}

%% Slides plus the presenter notes and the answers. `notes=show` puts beamer's
%% \note content on facing pages.
\DidactaDeclareProfile{slides-teacher}{beamer}{10pt,handout,notheorems}{%
  medium=slides,detail=brief,audience=teacher,solutions=full,pauses=off,notes=show}

%% ---------------------------------------------------------------------------
%% Document outputs
%% ---------------------------------------------------------------------------

%% The full written notes: everything on the slides, plus the prose and proofs
%% that do not fit on one. This is the student's book.
\DidactaDeclareProfile{notes}{article}{12pt,oneside}{%
  medium=document,detail=full,audience=student,solutions=hidden,pauses=off}

%% The same notes with worked solutions included.
\DidactaDeclareProfile{notes-solutions}{article}{12pt,oneside}{%
  medium=document,detail=full,audience=student,solutions=full,pauses=off}

%% The teacher's copy: full detail, teaching guidance, every solution.
\DidactaDeclareProfile{notes-teacher}{article}{12pt,oneside}{%
  medium=document,detail=full,audience=teacher,solutions=full,pauses=off}

%% A bound volume for a whole course: chapters, two-sided, real front matter.
\DidactaDeclareProfile{book}{book}{11pt,twoside,openright}{%
  medium=document,detail=full,audience=student,solutions=hidden,pauses=off}

%% ---------------------------------------------------------------------------
%% Single-sheet outputs
%% ---------------------------------------------------------------------------

%% A one- or two-page reference sheet, and the format a lab or seminar script
%% is built in. Tight margins, sections numbered but not subsections.
%%
%% Three levels, the same three a problem sheet has: a script with exercises in
%% it is handed out, checked and marked exactly like one.
\DidactaDeclareProfile{handout}{article}{11pt,oneside}{%
  medium=document,detail=full,audience=student,solutions=hidden,pauses=off,
  layout=compact}

%% The same sheet with the result under each exercise, so they can check
%% themselves.
\DidactaDeclareProfile{handout-answers}{article}{11pt,oneside}{%
  medium=document,detail=full,audience=student,solutions=answers,pauses=off,
  layout=compact}

%% The teacher's copy: results, worked solutions and the marking notes.
\DidactaDeclareProfile{handout-teacher}{article}{11pt,oneside}{%
  medium=document,detail=full,audience=teacher,solutions=full,pauses=off,
  layout=compact}

%% ---------------------------------------------------------------------------
%% Problem-sheet outputs
%%
%% The same source produces the three, and they are the three levels of the
%% problem editor: statement, result, worked solution. Level n shows fields 1
%% to n, which is what makes the menu answerable without reading a table.
%%
%% This is the axis the legacy system drove with \def\CurrentAudience and
%% multiaudience; here it is just `solutions`.
%%
%% There is deliberately no fourth `problems-solutions` -- a student copy with
%% the full working and no marking notes. It produced the same PDF as the
%% teacher's copy for every sheet in the repository (nothing uses `marking`),
%% so it was two menu entries for one output, and choosing between them was a
%% guess. The copy with everything in it is the teacher's.
%% ---------------------------------------------------------------------------

%% What the students receive: statements only.
\DidactaDeclareProfile{problems}{article}{12pt,oneside}{%
  medium=document,detail=full,audience=student,solutions=hidden,pauses=off,
  layout=compact}

%% Statements with the short result, so they can check their own work.
\DidactaDeclareProfile{problems-answers}{article}{12pt,oneside}{%
  medium=document,detail=full,audience=student,solutions=answers,pauses=off,
  layout=compact}

%% The teacher's copy: the worked solution step by step, the marking notes and
%% the common mistakes.
\DidactaDeclareProfile{problems-teacher}{article}{12pt,oneside}{%
  medium=document,detail=full,audience=teacher,solutions=full,pauses=off,
  layout=compact}

%% ---------------------------------------------------------------------------
%% Assessment output
%% ---------------------------------------------------------------------------

%% An exam paper. Same content, but numbered space for answers and no hints.
\DidactaDeclareProfile{exam}{article}{12pt,oneside}{%
  medium=document,detail=brief,audience=student,solutions=hidden,pauses=off,
  layout=exam}

%% The marking scheme for that exam.
\DidactaDeclareProfile{exam-marking}{article}{12pt,oneside}{%
  medium=document,detail=full,audience=teacher,solutions=full,pauses=off,
  layout=exam}

\didacta@profiles@restorecat


%% Fail loudly on an unknown profile. A silent fallback would hand back a PDF
%% that looks plausible and is the wrong output.
\@ifundefined{didacta@profile@\DidactaProfile @class}{%
  \begingroup
    \newlinechar=`\^^J
    \errmessage{^^J%
      ! Didacta: unknown output profile `\DidactaProfile'.^^J%
      ! Known profiles:\didacta@profile@list ^^J%
      ! Declare new ones with \string\DidactaDeclareProfile\space in
      didacta-profiles.tex.^^J}%
  \endgroup
}{}

\edef\didacta@class{\@nameuse{didacta@profile@\DidactaProfile @class}}
\edef\didacta@classoptions{\@nameuse{didacta@profile@\DidactaProfile @options}}

%% A document-level class override wins.
\ifx\DidactaClass\@empty\else
  \edef\didacta@class{\DidactaClass}%
\fi

%% Append the document's own options after the profile's, so the document wins
%% on any conflict.
\ifx\DidactaClassOptions\@empty\else
  \edef\didacta@classoptions{\didacta@classoptions,\DidactaClassOptions}%
\fi

%% `notes=show` has to reach beamer as a *class* option, not a later setting,
%% so it is the one axis the bootstrap has to look at. `\in@` is the kernel's
%% substring test, which keeps this dependency-free.
\edef\didacta@axes{\@nameuse{didacta@profile@\DidactaProfile @axes}}
\expandafter\in@\expandafter{\expandafter n\expandafter o\expandafter t%
  \expandafter e\expandafter s\expandafter =\expandafter s\expandafter h%
  \expandafter o\expandafter w\expandafter}\expandafter{\didacta@axes}
\ifin@
  \edef\didacta@classoptions{\didacta@classoptions,notes=show}%
\fi

\expandafter\documentclass\expandafter[\didacta@classoptions]{\didacta@class}

\makeatother

\usepackage{didacta}

%% Course metadata and the document title come from course.yaml and
%% year.yaml, in whichever language is being built. Didacta injects them.
%%
%% The declarations below are only a fallback, so that opening this file in an
%% editor and pressing compile still produces a titled document. Anything
%% Didacta injects overrides them.
\DidactaCourse{
  title   = {Anàlisi Matemàtica III},
  teacher = {Javier Falcó},
}
\DidactaDocument{Preliminars: espais normats}

\begin{document}
\DidactaTitlePage
\DidactaContents

\section{Espais normats}
\DidactaUnit{analysis/normed-spaces/definition}
\DidactaUnit{analysis/normed-spaces/induced-metric}

\end{document}
